\Needspace{20\baselineskip}\vspace{6mm}
\pdfbookmark[1]{The New Arrangements at the Mathematical and Physical Institute in Marburg.}{gerling-institut-1842}
\begin{center}\small Göttingische gelehrte Anzeigen.\quad 1842.\quad April 4.\end{center}

Professor \textsc{Gerling} in Marburg, a correspondent of the Royal Society of Sciences, wrote to Privy Councillor \textsc{Gauss} with a report on
\begin{center}\textit{the new arrangements at the mathematical and physical institute there.}\end{center}
We are all the more pleased to give an extract from that report here because the arrangements take account of the ordinary needs of teaching as well as of science itself, and the premises offer several distinctive features.

Since last autumn the institute has occupied a state building known as the Dörnberger Hof, which was remodelled and assigned to the university by order of His Highness the Electoral Prince and Co-Regent. In addition to the lecture hall, the second floor has four well-lit halls, a smaller room that can be darkened for optical experiments, and two separate workrooms. The instruments are set out with ample space so that each one can be used directly in its place without preventing or
\SourceLeaf{603}interfering with the simultaneous use of others, and so that students who are so inclined and sufficiently prepared can carry out exercises of their own conveniently. The institute already has a substantial stock of physical apparatus, which is further completed and enlarged every year. With the connecting doors open, these rooms provide an unobstructed straight line 150 feet long, an advantage of real value for a variety of purposes.

The director’s official residence is on the third floor. The institute also has its own workshop and forge, which have been placed on the ground floor.

For experiments that must be carried out in the open air, there is ample opportunity both around the building and on the platform of a tower founded on the rock and directly connected to it. The tower also offers an extremely valuable resource for certain kinds of physical work, namely an unobstructed drop of about 80 feet. To obtain this, square openings have been made in the floors of the three storeys of the tower and in the floor above the cellar (formerly a dungeon). On the upper platform, which is reached from the roof of the building by a spiral staircase, there is also an octagonal pavilion 15 feet in diameter, with a similar opening in its floor. which can be covered with a trapdoor or surrounded by a gallery, as circumstances require. The tower also contains a very solid wall nearly 20 feet high, which may offer important advantages on some occasions.

These arrangements also met a need that is felt particularly at a university with no observatory. Until now, for example, the time determinations that have become indispensable for so many kinds of physical work could there be obtained only by the time-consuming method of corresponding altitudes of the Sun. All that was needed was to arrange the doors and windows of the pavilion appropriately, so that all the required observations could be made there with the movable astronomical instruments, some of which the institute had long possessed. An \textsc{Ertel} portable transit instrument, for example, has a completely unobstructed range from its usual station, from horizon to horizon in the meridian, and from the western horizon to an easterly zenith distance of about 11 degrees in the prime vertical. The position of this site
\SourceLeaf{604}was determined by Professor \textsc{Gerling}’s trigonometrical survey of the Electorate, linked to the Hanoverian arc survey. We look forward eagerly to the publication of its details. The latitude is $50^0\;48'\;46''9$, and the longitude from Ferro $26^0\;26'\;2''3$.

Finally, one more arrangement deserves mention. It was considered essential to completing the whole plan. However exceptionally well the main building suited all the other purposes, it had the drawback that it was least suited, relatively speaking, to meteorological observations, since the considerably higher Schlossberg towers over it. Besides, for various purposes it remains very desirable to have two sites quite far apart yet accessible from each other. Both needs were met by converting an old turret on the highest point of the Schlossberg, 1900 feet away in a straight line and about 100 feet higher, into a meteorological tower. This made it possible, among other things, to set up a lightning conductor for observations away from the town and thus to eliminate every conceivable danger. It also made it possible, for example, to carry out experiments there in the near future relating to magnetic telegraphy.

